% !TeX program = xelatex
% !TeX encoding = UTF-8
% ===========================================================================
% 全国大学生电子设计竞赛 — 控制类通用技术报告模板
% 编译：xelatex template_general.tex （建议运行两次以更新目录和引用）
% ===========================================================================
\documentclass[12pt,a4paper,fontset=fandol]{ctexart}

% ---------- 页面设置（按竞赛要求） ----------
\usepackage[top=3.5cm,bottom=2.5cm,left=2.5cm,right=2.5cm]{geometry}
%   ↑ top=3.5cm 保证上方留出 ≥3cm 空白（含页眉 0.5cm 余量）
\usepackage{setspace}
\setstretch{1.25}

% ---------- 数学与符号 ----------
\usepackage{amsmath}
\usepackage{amssymb}

% ---------- 表格 ----------
\usepackage{booktabs}
\usepackage{tabularx}
\usepackage{longtable}
\usepackage{array}

% ---------- 图片 ----------
\usepackage{graphicx}
\usepackage{float}

% ---------- 题注 ----------
\usepackage[font=small]{caption}
\captionsetup{labelsep=space}

% ---------- 列表 ----------
\usepackage{enumitem}
\setlist{nosep}

% ---------- 页眉页脚（页码右下） ----------
\usepackage{fancyhdr}
\pagestyle{fancy}
\fancyhf{}
\fancyfoot[R]{\thepage}       % 页码在右下
\renewcommand{\headrulewidth}{0pt}

% ---------- 超链接 ----------
\usepackage[hidelinks]{hyperref}

% ---------- 代码 ----------
\usepackage{listings}
\lstset{
  basicstyle=\ttfamily\small,
  breaklines=true,
  frame=single,
  numbers=left,
  numberstyle=\tiny,
  showstringspaces=false,
}

% ---------- 标题格式 ----------
\ctexset{
  section      = {format = \large\bfseries},
  subsection   = {format = \normalsize\bfseries},
  subsubsection = {format = \normalsize\bfseries},
}

% ---------- 工具命令 ----------
\newcolumntype{Y}{>{\raggedright\arraybackslash}X}
\newcommand{\code}[1]{\texttt{\small #1}}
\newcommand{\placeholder}[1]{\textcolor{red}{\texttt{<#1>}}}
%   ↑ 红色占位符，填写后改为正常文字

\usepackage{xcolor}  % 用于 \placeholder 红色占位符高亮

\begin{document}

% =========================== 封 面 ===========================
\begin{titlepage}
  \centering
  \vspace*{2.2cm}
  {\zihao{2}\bfseries \placeholder{年份}年全国大学生电子设计竞赛\par}
  \vspace{0.5cm}
  {\zihao{2}\bfseries \placeholder{赛区}赛区（\placeholder{赞助商}杯）\par}
  \vfill
  {\zihao{1}\bfseries \placeholder{作品名称}\par}
  \vspace{0.7cm}
  {\zihao{3}\bfseries 题目编号：\placeholder{题号}\par}
  \vfill
  \begin{tabular}{rl}
    \zihao{4}参赛队编号：& \rule{7cm}{0.4pt}\\[0.55cm]
    \zihao{4}参赛队学校：& \rule{7cm}{0.4pt}\\[0.55cm]
    \zihao{4}参赛队学生：& \rule{7cm}{0.4pt}
  \end{tabular}
  \vfill
  {\zihao{4}\placeholder{二零XX年X月}\par}
\end{titlepage}

% =========================== 正 文 ===========================
\setcounter{page}{1}
\begin{center}
  {\zihao{2}\bfseries \placeholder{作品名称}}
\end{center}

% ---------- 摘 要 ----------
\begin{abstract}
\noindent
% ===== 以下为摘要正文（≤300字），请替换 =====
针对\placeholder{题目名称}中\placeholder{核心任务描述}的需求，设计了一套以\placeholder{主控芯片型号}为主控的\placeholder{系统名称}。
系统采用\placeholder{传感器1}、\placeholder{传感器2}和\placeholder{执行器}构成闭环控制系统。
控制层采用\placeholder{控制算法名称}实现\placeholder{核心功能}；
任务层采用\placeholder{状态机/调度策略}将\placeholder{子任务列表}统一调度。

系统实现了\placeholder{基本要求1}、\placeholder{基本要求2}和\placeholder{发挥部分1}，
\placeholder{量化指标：如误差\textless X\%, 用时\textless Y s, 成功率达Z\%}。
本文给出硬件组成、关键控制公式、软件流程和可复现的测试方案。

\noindent\textbf{关键词：}\placeholder{关键词1}；\placeholder{关键词2}；\placeholder{关键词3}；\placeholder{关键词4}；\placeholder{关键词5}
\end{abstract}

% =========================== 第 1 章 ===========================
\section{设计方案与工作原理}

\subsection{预期实现目标定位}
% === 替换为你的题目要求分析 ===
本设计面向\placeholder{题号}题规定的任务：
\begin{itemize}
  \item 基本要求（1）：\placeholder{描述基本要求1}
  \item 基本要求（2）：\placeholder{描述基本要求2}
  \item 发挥部分（1）：\placeholder{描述发挥部分1}
  \item 发挥部分（2）：\placeholder{描述发挥部分2}
\end{itemize}
设计须遵守以下限制条件：
\begin{itemize}
  \item \placeholder{限制条件1，如：必须使用 TI MSPM0 系列 MCU}
  \item \placeholder{限制条件2，如：不得安装摄像头}
  \item \placeholder{限制条件3，如：车载电池供电，不可外接电源/设备}
\end{itemize}

\subsection{技术方案分析比较}
% === 每个模块至少给出 2~3 个方案，说明各自优劣及最终选择理由 ===

\textbf{1. \placeholder{核心功能}方案论证。}
方案一：\placeholder{方案一描述}，
其优点为\placeholder{优点}，缺点为\placeholder{缺点}。
方案二：\placeholder{方案二描述}，
其优点为\placeholder{优点}，缺点为\placeholder{缺点}。
方案三：\placeholder{方案三描述}，
其优点为\placeholder{优点}，且\placeholder{补充优势}。
\emph{方案选择：}方案三兼顾\placeholder{性能指标}与\placeholder{实现复杂度}，
因此\placeholder{最终选择方案三}。

\textbf{2. \placeholder{传感器/检测}方案论证。}
方案一：\placeholder{方案一描述}。
方案二：\placeholder{方案二描述}。
方案三：\placeholder{方案三描述}。
\emph{方案选择：}\placeholder{从精度、可靠性、成本等方面论述选择理由}。

\textbf{3. \placeholder{控制算法}方案论证。}
方案一：\placeholder{如开环控制}，\placeholder{优缺点}。
方案二：\placeholder{如 PID 控制}，\placeholder{优缺点}。
方案三：\placeholder{如模糊/自适应/串级 PID}，\placeholder{优缺点}。
\emph{方案选择：}\placeholder{结合题目要求论述选择理由}。

\subsection{系统结构工作原理}
% === 给出系统总体框图描述 ===
系统由\placeholder{感知、控制、执行、人机交互}四部分组成。
\placeholder{传感器部分}负责\placeholder{数据采集}；
\placeholder{主控芯片}在\placeholder{控制周期}定时中断内完成采样、控制算法运算与输出刷新；
\placeholder{执行器}通过\placeholder{驱动方式}驱动\placeholder{执行机构}；
\placeholder{显示/提示模块}显示运行状态。

\begin{table}[H]
\centering
\caption{系统模块组成}
\begin{tabularx}{\textwidth}{p{2.2cm}p{4cm}Y}
\toprule
模块 & 主要器件/接口 & 作用\\
\midrule
主控 & \placeholder{MCU 型号} & \placeholder{功能描述}\\
执行 & \placeholder{驱动 + 执行器} & \placeholder{功能描述}\\
反馈 & \placeholder{传感器} & \placeholder{功能描述}\\
人机交互 & \placeholder{显示/按键/指示灯} & \placeholder{功能描述}\\
\bottomrule
\end{tabularx}
\end{table}

% 如有场地/系统示意图，取消下面注释
% \begin{figure}[H]
% \centering
% \includegraphics[width=0.85\textwidth]{figures/system_diagram.png}
% \caption{系统总体框图 / 场地示意图}
% \end{figure}

\subsection{功能指标实现方法}
% === 分条说明每个功能/指标如何实现 ===
\begin{enumerate}
  \item \textbf{\placeholder{功能1}：}\placeholder{实现方法概述，含关键算法和处理流程}
  \item \textbf{\placeholder{功能2}：}\placeholder{实现方法概述}
  \item \textbf{\placeholder{功能3}：}\placeholder{实现方法概述}
\end{enumerate}

\subsection{测量、控制与分析处理}
% === 描述控制周期的时序和数据流 ===
每个\placeholder{控制周期}控制周期中，
首先\placeholder{读取传感器数据并预处理}，
随后\placeholder{运行控制算法计算输出量}，
最后由\placeholder{执行器输出}。
\placeholder{其他传感器}数据以\placeholder{更新频率}节拍更新，
用于\placeholder{辅助功能}。

% =========================== 第 2 章 ===========================
\section{核心部件与电路设计}

\subsection{关键器件性能分析}
% === 逐一说明关键器件选型依据和性能参数 ===
\begin{itemize}
  \item \textbf{\placeholder{主控芯片}：}\placeholder{核心参数：内核、主频、外设资源等}，满足\placeholder{控制频率、外设数量}需求。
  \item \textbf{\placeholder{传感器1}：}\placeholder{关键指标：量程、精度、分辨率、接口}。
  \item \textbf{\placeholder{执行器/驱动器}：}\placeholder{关键指标：电压、电流、接口方式}。
\end{itemize}

\subsection{电路结构工作机理}
% === 说明关键电路连接和控制方式 ===
\placeholder{描述关键电路接口和控制信号流向。如：PWM 输出引脚→电机驱动→电机；
传感器信号→ADC/GPIO/UART→MCU 处理。}

\subsection{运动学/控制模型}
% === 给出核心数学公式（LaTeX 格式） ===
设\placeholder{系统状态变量定义}，则系统模型为：
\begin{equation}
  \placeholder{核心公式1}
  \label{eq:key1}
\end{equation}
\begin{equation}
  \placeholder{核心公式2}
  \label{eq:key2}
\end{equation}

控制器以\placeholder{误差定义}为输入：
\begin{equation}
  u(k) = \placeholder{控制律表达式}
  \label{eq:controller}
\end{equation}
式中，\placeholder{参数说明}。

\subsection{关键参数与接口}
\begin{table}[H]
\centering
\caption{控制系统关键参数}
\begin{tabularx}{\textwidth}{p{3.2cm}p{3.6cm}Y}
\toprule
参数 & 取值 & 说明\\
\midrule
控制周期 & \placeholder{如 5\,ms} & \placeholder{说明}\\
\placeholder{参数1} & \placeholder{取值} & \placeholder{说明}\\
\placeholder{参数2} & \placeholder{取值} & \placeholder{说明}\\
\placeholder{参数3} & \placeholder{取值} & \placeholder{说明}\\
\bottomrule
\end{tabularx}
\end{table}

% =========================== 第 3 章 ===========================
\section{系统软件设计分析}

\subsection{系统总体工作流程}
% === 描述上电→初始化的流程，以及主循环/中断的职责 ===
上电后完成\placeholder{外设初始化}和\placeholder{参数配置}。
系统采用\placeholder{定时中断+主循环}架构：
\begin{itemize}
  \item \textbf{定时中断（\placeholder{频率}）：}\placeholder{传感器采样→控制算法→执行器输出}
  \item \textbf{主循环：}\placeholder{显示刷新、串口通信、按键处理、数据记录}
\end{itemize}

\subsection{任务状态机设计}
% === 如果系统有多个状态/阶段，用状态机描述；否则可删除此节 ===
系统将\placeholder{任务名称}分解为以下状态：
\begin{table}[H]
\centering
\caption{任务状态机}
\begin{tabularx}{\textwidth}{p{2.6cm}p{2.5cm}Yp{2.7cm}}
\toprule
状态 & 输入依据 & 主要动作 & 退出条件\\
\midrule
\code{\placeholder{STATE\_1}} & \placeholder{传感器数据} & \placeholder{执行动作} & \placeholder{切换判据}\\
\code{\placeholder{STATE\_2}} & \placeholder{传感器数据} & \placeholder{执行动作} & \placeholder{切换判据}\\
\code{\placeholder{STATE\_3}} & \placeholder{传感器数据} & \placeholder{执行动作} & \placeholder{切换判据}\\
\code{STOP} & — & 输出清零、声光提示 & 任务完成\\
\bottomrule
\end{tabularx}
\end{table}

\subsection{关键模块程序设计}
% === 列出核心函数/模块的功能说明 ===
\begin{itemize}
  \item \code{\placeholder{Module1\_Init()}}：\placeholder{初始化说明}
  \item \code{\placeholder{Module1\_Process()}}：\placeholder{处理逻辑说明}
  \item \code{\placeholder{Control\_Loop()}}：\placeholder{控制循环说明}
  \item \code{\placeholder{Sensor\_Read()}}：\placeholder{传感器读取说明}
\end{itemize}

\subsection{可靠性处理}
% === 异常检测、容错、抗干扰措施 ===
为保障系统可靠性，采取了以下措施：
\begin{enumerate}
  \item \textbf{传感器异常检测：}\placeholder{如超时检测、数据范围校验、通信校验}
  \item \textbf{状态切换防抖：}\placeholder{连续 N 个周期条件满足才切换状态}
  \item \textbf{输出限幅保护：}\placeholder{PWM/电压/电流限幅范围}
  \item \textbf{超时/安全停止：}\placeholder{如最大运行时间保护、手动急停}
\end{enumerate}

% =========================== 第 4 章 ===========================
\section{竞赛工作环境条件}

\subsection{软件环境}
\begin{itemize}
  \item IDE / 编译工具：\placeholder{如 Keil uVision5、IAR、CCS}
  \item 目标工程：\code{\placeholder{工程名称}}
  \item 外设配置工具：\placeholder{如 TI SysConfig、STM32CubeMX}
  \item 调试工具：\placeholder{如串口助手、VOFA+、DataScope、逻辑分析仪}
\end{itemize}

\subsection{硬件平台与场地}
硬件平台包括：
\begin{itemize}
  \item \placeholder{主控板型号}
  \item \placeholder{传感器列表及型号}
  \item \placeholder{执行器列表及型号}
  \item \placeholder{供电方案}
\end{itemize}
测试场地：\placeholder{场地尺寸和制作要求描述}。

\subsection{调试条件与安全事项}
\begin{enumerate}
  \item 先在\placeholder{架空/空载}状态确认\placeholder{输出极性、方向、急停逻辑}；
  \item 再在\placeholder{低速/低功率}下完成\placeholder{参数整定}；
  \item \placeholder{供电安全注意事项}；
  \item \placeholder{机械/运动安全注意事项}。
\end{enumerate}

% =========================== 第 5 章 ===========================
\section{作品成效总结分析}

\subsection{测试方案}
% === 从简到难设计测试步骤 ===
按由简到难的顺序进行测试：
\begin{enumerate}
  \item \textbf{单元测试：}\placeholder{各模块独立功能验证}
  \item \textbf{集成测试：}\placeholder{基本要求逐项验证}
  \item \textbf{系统测试：}\placeholder{发挥部分和极限条件验证}
  \item \textbf{重复性测试：}\placeholder{每组条件至少重复 3 次}
\end{enumerate}
测试仪器：\placeholder{具体型号，如 Tektronix MDO3024 示波器、FLUKE 17B+ 万用表}。

\begin{table}[H]
\centering
\caption{实车/实物测试记录表}
\small
\begin{tabularx}{\textwidth}{p{2.3cm}Yp{1.2cm}p{1.2cm}p{1.2cm}p{1.5cm}}
\toprule
测试项目 & 评价指标 & 第1次 & 第2次 & 第3次 & 结论\\
\midrule
\placeholder{基本要求(1)} & \placeholder{指标} & \placeholder{待实测} & \placeholder{待实测} & \placeholder{待实测} & \placeholder{待判定}\\
\placeholder{基本要求(2)} & \placeholder{指标} & \placeholder{待实测} & \placeholder{待实测} & \placeholder{待实测} & \placeholder{待判定}\\
\placeholder{发挥部分(1)} & \placeholder{指标} & \placeholder{待实测} & \placeholder{待实测} & \placeholder{待实测} & \placeholder{待判定}\\
\placeholder{发挥部分(2)} & \placeholder{指标} & \placeholder{待实测} & \placeholder{待实测} & \placeholder{待实测} & \placeholder{待判定}\\
\bottomrule
\end{tabularx}
\end{table}

\subsection{成效与不足分析}
% === 实事求是：讲清楚做到了什么、还有什么问题 ===
\subsubsection*{系统成效}
\placeholder{逐项对照题目要求，说明达标情况。使用量化数据支撑结论。}

\subsubsection*{存在不足}
\placeholder{分析存在的问题：如精度受限、抗干扰不足、特定场景表现不稳定等，说明原因。}

\subsection{创新特色与展望}
\subsubsection*{创新特色}
\placeholder{总结本设计的创新点，如：算法改进、结构优化、低成本高精度方案、多传感器融合策略等。}

\subsubsection*{改进方向}
\placeholder{提出后续可优化的方向，如：自适应参数、机器学习、更高精度传感器等。}

% =========================== 第 6 章 ===========================
\section{参考资料及文献}
\begin{enumerate}[label={[\arabic*]}]
  \item \placeholder{年份}年全国大学生电子设计竞赛\placeholder{赛区}赛区\placeholder{题号}题试题。
  \item \placeholder{芯片/传感器数据手册}。
  \item \placeholder{参考书籍、论文或在线资源}。
  \item \placeholder{项目内部文档引用}。
\end{enumerate}

% =========================== 第 7 章 ===========================
\section{附件材料}

\subsection{工程文件与关键源代码}
工程文件位于\code{\placeholder{固件路径/}}。
核心源代码文件：
\begin{itemize}
  \item \code{\placeholder{control.c}}：\placeholder{功能说明}
  \item \code{\placeholder{sensor.c}}：\placeholder{功能说明}
  \item \code{\placeholder{motor.c}}：\placeholder{功能说明}
\end{itemize}

\subsection{核心源代码清单}
% === 粘贴关键代码片段（建议不超过 2 页） ===
%\begin{lstlisting}[language=C, caption={\placeholder{核心控制函数}}]
%// 在此粘贴关键控制代码
%\end{lstlisting}

\subsection{最终提交前检查清单}
\begin{itemize}
  \item[$\square$] \textbf{格式核验：}封面和正文无学校、代码、姓名；页面上方空白 ≥3cm；A4 纸 6~8 页。
  \item[$\square$] \textbf{内容核验：}测试仪器均已标注具体型号；每个方案均有 ≥2 个对比方案；公式按 LaTeX 格式。
  \item[$\square$] \textbf{数据核验：}所有实测数据已补录并复核；无虚构数据。
  \item[$\square$] \textbf{规则核验：}芯片/器件选型符合题目限制；无禁用的传感器或外部设备。
  \item[$\square$] \textbf{完整性核验：}摘要、正文、参考文献、附件齐全；页码连续。
\end{itemize}

\end{document}
